\documentclass[12pt,a4paper]{article}
\usepackage[UTF8]{ctex}
\usepackage{amsmath,amssymb,amsthm}
\usepackage{geometry}
\usepackage{bm}
\usepackage{hyperref}

\geometry{margin=2.5cm}

\title{力和运动自由度的对偶性详解}
\author{现代机器人学：第12章 抓取与操作\\
详细解释}
\date{}

\begin{document}

\maketitle

\section{核心观点}

文本中的关键陈述是：

\begin{quote}
对于任何点接触和接触标签，由该接触引起的物体运动的等式约束数量等于它提供的wrench自由度数量。
\end{quote}

这是一个深刻的\textbf{对偶性原理}，连接了运动学和静力学。本节将详细解释这个原理。

\section{基本概念}

\subsection{运动约束（等式约束）}

\textbf{运动约束}是指限制物体运动的等式方程。对于刚体，这些约束限制了twist $V$的某些分量。

\begin{itemize}
    \item \textbf{约束数量}：等于独立等式约束的数量
    \item \textbf{物理意义}：物体必须满足这些条件才能保持特定的接触模式
    \item \textbf{数学形式}：通常是$F^T V = 0$或$\dot{p}_A = \dot{p}_B$等形式
\end{itemize}

\subsection{Wrench自由度}

\textbf{Wrench自由度}是指可以通过接触传递的独立wrench的数量或维度。

\begin{itemize}
    \item \textbf{物理意义}：接触可以施加多少种不同的力/力矩组合
    \item \textbf{数学形式}：wrench锥的维度
    \item \textbf{约束关系}：wrench必须位于摩擦锥内，且与接触模式一致
\end{itemize}

\subsection{对偶性的直观理解}

\begin{center}
\textbf{约束运动越多 $\Leftrightarrow$ 能施加的力越少}
\end{center}

\begin{center}
\textbf{约束运动越少 $\Leftrightarrow$ 能施加的力越多}
\end{center}

这是一个自然的平衡：如果你限制了物体的某些运动，那么接触就可以在这些受限方向上施加力；如果你允许更多运动，那么接触在这些方向上的力传递能力就受限。

\section{三维情况分析}

\subsection{分离接触（B：Breaking）}

\textbf{定义}：接触正在分离，满足不可穿透约束但不满足主动约束。

\subsubsection{运动约束}

\begin{itemize}
    \item \textbf{约束数量}：0个等式约束
    \item \textbf{原因}：物体可以自由运动，只要不穿透（不等式约束$\dot{d} \geq 0$）
    \item \textbf{数学表达}：没有等式约束，只有不等式$F^T V \geq 0$
\end{itemize}

\subsubsection{Wrench自由度}

\begin{itemize}
    \item \textbf{自由度数量}：0
    \item \textbf{原因}：接触正在分离，无法传递力
    \item \textbf{物理意义}：$f_n = 0$（法向力为零），因此无法施加任何接触力
\end{itemize}

\textbf{结论}：$0$个运动约束 = $0$个wrench自由度 ✓

\subsection{滑动接触（S：Sliding）}

\textbf{定义}：接触保持（满足主动约束$F^T V = 0$），但存在相对切向运动。

\subsubsection{运动约束}

\begin{itemize}
    \item \textbf{约束数量}：1个等式约束
    \item \textbf{约束方程}：$F^T (V_A - V_B) = 0$（或$F^T V_A = 0$，如果$B$静止）
    \item \textbf{物理意义}：必须满足一个方程以保持接触（法向速度为零）
    \item \textbf{数学表达}：$\hat{n}^T (\dot{p}_A - \dot{p}_B) = 0$
\end{itemize}

这个约束限制了twist空间的一个维度。对于三维物体（6维twist空间），滑动接触将可行twist限制在5维超平面上。

\subsubsection{Wrench自由度}

\begin{itemize}
    \item \textbf{自由度数量}：1
    \item \textbf{原因}：对于给定的滑动运动，接触wrench只有一个自由度
    \item \textbf{物理意义}：接触力必须位于摩擦锥的边缘上，且方向与滑动方向相反
    \item \textbf{数学表达}：$F = k F_{\text{edge}}$，其中$k \geq 0$是标量，$F_{\text{edge}}$是摩擦锥边缘上的wrench
\end{itemize}

具体来说：
\begin{itemize}
    \item 法向力$f_n$的大小可以变化（一个自由度）
    \item 切向力$f_t$的方向和大小由滑动方向决定：$f_t = \mu f_n$，方向与滑动方向相反
    \item 因此，整个wrench由法向力大小$f_n$唯一确定（一个标量参数）
\end{itemize}

\textbf{结论}：$1$个运动约束 = $1$个wrench自由度 ✓

\subsection{滚动/固定接触（R：Rolling/Sticking）}

\textbf{定义}：接触保持，且接触点处没有相对运动（$\dot{p}_A = \dot{p}_B$）。

\subsubsection{运动约束}

\begin{itemize}
    \item \textbf{约束数量}：3个等式约束
    \item \textbf{约束方程}：$\dot{p}_A = \dot{p}_B$，即
    \begin{equation}
    v_A + [\omega_A]p_A = v_B + [\omega_B]p_B
    \end{equation}
    \item \textbf{展开形式}：
    \begin{align}
    v_{Ax} - \omega_{Az} p_{Ay} + \omega_{Ay} p_{Az} &= v_{Bx} - \omega_{Bz} p_{By} + \omega_{By} p_{Bz} \\
    v_{Ay} + \omega_{Az} p_{Ax} - \omega_{Ax} p_{Az} &= v_{By} + \omega_{Bz} p_{Bx} - \omega_{Bx} p_{Bz} \\
    v_{Az} + \omega_{Ax} p_{Ay} - \omega_{Ay} p_{Ax} &= v_{Bz} + \omega_{Bx} p_{By} - \omega_{By} p_{Bx}
    \end{align}
    \item \textbf{物理意义}：物体上接触点的运动被完全指定（三个方向的速度分量）
    \item \textbf{数学表达}：对于6维twist空间，滚动约束将可行twist限制在3维子空间中
\end{itemize}

\subsubsection{Wrench自由度}

\begin{itemize}
    \item \textbf{自由度数量}：3
    \item \textbf{原因}：接触wrench锥内部的任何wrench都与接触模式一致
    \item \textbf{物理意义}：
    \begin{itemize}
        \item 法向力$f_n$可以变化（1个自由度）
        \item 切向力$f_t$可以在摩擦锥内任意方向，大小满足$|f_t| \leq \mu f_n$（2个自由度：方向和大小）
    \end{itemize}
    \item \textbf{数学表达}：wrench锥$\mathcal{WC} = \text{pos}(\{F_1, F_2, F_3, \ldots\})$是3维的
\end{itemize}

对于无摩擦接触，wrench自由度是1（只有法向力）。对于有摩擦接触，如果摩擦锥是3维的（三维空间中的锥），则wrench自由度是3。

\textbf{注意}：文本中提到"接触wrench锥内部的任何wrench"，这意味着wrench可以在整个摩擦锥内自由选择，因此有3个自由度（法向力大小 + 切向力的2个方向分量）。

\textbf{结论}：$3$个运动约束 = $3$个wrench自由度 ✓

\section{平面情况分析}

对于平面问题，twist空间是3维的（$\omega_z, v_x, v_y$），wrench空间也是3维的（$m_z, f_x, f_y$）。

\subsection{分离接触（B）}

\begin{itemize}
    \item \textbf{运动约束}：0个
    \item \textbf{Wrench自由度}：0个
\end{itemize}

\textbf{结论}：$0 = 0$ ✓

\subsection{滑动接触（S）}

\begin{itemize}
    \item \textbf{运动约束}：1个（$F^T V = 0$，其中$F = [m_z, f_x, f_y]^T$）
    \item \textbf{Wrench自由度}：1个（摩擦锥边缘上的wrench，由法向力大小确定）
\end{itemize}

\textbf{结论}：$1 = 1$ ✓

\subsection{滚动/固定接触（R）}

\begin{itemize}
    \item \textbf{运动约束}：2个
    \begin{itemize}
        \item 约束1：$F^T V = 0$（保持接触）
        \item 约束2：切向速度为零（$\dot{p}_{A,\text{tan}} = \dot{p}_{B,\text{tan}}$）
        \item 注意：在平面中，切向方向是1维的，所以只需要1个额外的约束
        \item 总共：1（法向）+ 1（切向）= 2个约束
    \end{itemize}
    \item \textbf{Wrench自由度}：2个
    \begin{itemize}
        \item 平面摩擦锥是2维的（由两条边正张成）
        \item 法向力大小 + 切向力方向/大小 = 2个自由度
    \end{itemize}
\end{itemize}

\textbf{结论}：$2 = 2$ ✓

\section{对偶性的数学基础}

\subsection{约束空间与力空间的对偶}

在刚体力学中，存在一个基本的对偶关系：

\begin{center}
\textbf{约束空间} $\Leftrightarrow$ \textbf{力空间}
\end{center}

\begin{itemize}
    \item 每个运动约束对应一个可以施加力的方向
    \item 约束的维度 = 可以施加的力的维度
\end{itemize}

\subsection{虚功原理}

对偶性可以通过虚功原理理解：

\begin{equation}
F^T V = 0
\end{equation}

其中：
\begin{itemize}
    \item $F$是wrench（力/力矩）
    \item $V$是twist（速度/角速度）
    \item 如果$V$被约束在某个子空间中，则$F$可以在对偶子空间中自由变化
\end{itemize}

\subsection{线性代数的视角}

\begin{itemize}
    \item 运动约束定义了twist空间的一个子空间（约束子空间）
    \item Wrench自由度对应wrench空间中与约束子空间对偶的子空间
    \item 根据线性代数，约束子空间的维度 + 对偶子空间的维度 = 总空间维度
    \item 但在这个对偶性中，约束维度 = 对偶维度（wrench自由度）
\end{itemize}

\section{实际应用意义}

\subsection{接触模式识别}

理解对偶性有助于：
\begin{itemize}
    \item 快速判断接触模式：如果知道运动约束数量，立即知道wrench自由度
    \item 验证分析正确性：检查约束数量和自由度是否匹配
\end{itemize}

\subsection{抓取规划}

在抓取规划中：
\begin{itemize}
    \item 如果需要更多约束（形状闭合），需要更多接触
    \item 如果需要传递更多力（力闭合），需要合适的接触配置
    \item 对偶性帮助理解约束和力传递之间的平衡
\end{itemize}

\subsection{接触力分析}

在分析接触力时：
\begin{itemize}
    \item 知道运动约束后，可以确定有多少个独立的力可以施加
    \item 帮助理解为什么某些接触模式可以传递某些力而不能传递其他力
\end{itemize}

\section{总结}

\begin{enumerate}
    \item \textbf{核心原理}：运动约束数量 = Wrench自由度数量
    
    \item \textbf{三维情况}：
    \begin{itemize}
        \item B（分离）：0约束 = 0自由度
        \item S（滑动）：1约束 = 1自由度
        \item R（滚动）：3约束 = 3自由度
    \end{itemize}
    
    \item \textbf{平面情况}：
    \begin{itemize}
        \item B（分离）：0约束 = 0自由度
        \item S（滑动）：1约束 = 1自由度
        \item R（滚动）：2约束 = 2自由度
    \end{itemize}
    
    \item \textbf{物理直觉}：约束越多，能施加的力越受限；约束越少，能施加的力越多
    
    \item \textbf{数学基础}：基于虚功原理和线性代数的对偶性
\end{enumerate}

这个对偶性是接触力学中的一个基本原理，连接了运动学和静力学，为理解和分析接触问题提供了强大的工具。

\end{document}

